\begin{pastexam}[label={px:gk3-102-ds-06}]{高考三級資訊處理 102 年 第六題}
（申論；6 分）

有一雜湊表格（hash table）T 的記憶空間共含 11 個桶（buckets），位址編號由 0 至 10，每個桶有一個槽
（slot）。雜湊函數 h1 定義為 h1(key) = key \% 11，當有碰撞（collision）發生時採二次雜湊開放定址法
（open addressing with double hashing）處理，其函數定義為 h(key, j) = (h1(key)＋j * h2(key)) \% 11，
其中 j 為碰撞次數，j = 1, 2, 3, ..., 11，h2(key) = 1＋(key \% 10)。欲將 26 放入雜湊表格 T，總共經過
6 次探測才成功找到存放位址。請問 26 在雜湊表格 T 的探測順序為何？（6 分）
\end{pastexam}

\begin{pastsoln}
答案：探測的位址依序是 $4, 0, 7, 3, 10, 6$，26 最後放在位址 $6$。

第一次探測是 $h_1(26)$，之後第 $j$ 次碰撞後探測 $h(26, j)$（資結書第~12.3~節的雙重雜湊）：
\[
h_1(26) = 26 \bmod 11 = 4, \qquad h_2(26) = 1 + (26 \bmod 10) = 7 .
\]
\begin{center}
\small
\begin{tabular}{ccl}
\toprule
第幾次探測 & $j$ & 位址 \\
\midrule
1 & 0 & $h_1(26) = 4$ \\
2 & 1 & $(4 + 1 \cdot 7) \bmod 11 = 11 \bmod 11 = 0$ \\
3 & 2 & $(4 + 2 \cdot 7) \bmod 11 = 18 \bmod 11 = 7$ \\
4 & 3 & $(4 + 3 \cdot 7) \bmod 11 = 25 \bmod 11 = 3$ \\
5 & 4 & $(4 + 4 \cdot 7) \bmod 11 = 32 \bmod 11 = 10$ \\
6 & 5 & $(4 + 5 \cdot 7) \bmod 11 = 39 \bmod 11 = 6$ \\
\bottomrule
\end{tabular}
\end{center}
「經過 6 次探測才成功」表示前五個位址 $4, 0, 7, 3, 10$ 都已經有資料（碰撞了 $5$ 次），第六次的
位址 $6$ 是空的，26 放在那裡。 實際上每一步只要把前一個位址加 $7$ 再對 $11$ 取餘數就好。

順帶一提，桶數 $11$ 是質數，$h_2$ 的值 $1$ 到 $10$ 都和 $11$ 互質，所以探測序列一定會走遍全部 $11$ 個
位址（資結書定理~12.4），只要表沒滿就一定找得到空位。 題目寫 $j = 1, \ldots, 11$，其實 $j = 11$
時 $(4 + 77) \bmod 11 = 4$ 又回到起點，有用的只到 $j = 10$。
\end{pastsoln}
